\documentclass[10pt,a4paper]{amsart}
% Este arquivo sera processado automaticamente por um programa.
% Por favor não altere o preâmbulo (antes de \begin{document} ).

% This file will be compiled by a computer program.
% Please do not change the preamble (before \begin{document} ).

\usepackage{amsmath,amssymb}
\usepackage{enumerate}
\usepackage{epsfig,geometry,graphicx,graphics,enumitem,amsthm,indentfirst}
\usepackage[english]{babel}

% Por favor nao inclua outros "packages".
% Please do not include any other package.

\begin{document}


\title{The Mapper algorithm: a discrete analogue of the Reeb graph}
\author{Guilherme Vituri}


\maketitle

The Reeb graph of a continuous function defined on a topological space is a
classical object that records how the connected components of its level sets
evolve, yielding a compact one-dimensional summary of the space~\cite{reeb1946}.
When the space is known only through a finite sample, this construction cannot
be applied directly. The Mapper algorithm, introduced by Singh, M\'emoli and
Carlsson~\cite{singh2007}, overcomes this limitation: given a point cloud, a
filter function, and a cover of its image, Mapper clusters the preimage of each
cover element and assembles the resulting clusters into the nerve of the refined
cover, producing a graph. This graph is a discrete analogue of the Reeb graph,
and under suitable hypotheses it is provably close to it~\cite{carriere2018}. In
this talk we present the Mapper algorithm from this topological perspective,
discuss its main variants and the stability results that justify it, and
illustrate the constructions computationally with the \texttt{TDAmapper.jl}
package, closing with an application to exploratory data
analysis~\cite{lum2013}.
\\

	\begin{thebibliography}{99}

		\bibitem{reeb1946} Reeb, G. Sur les points singuliers d'une forme de
		Pfaff compl\`etement int\'egrable ou d'une fonction num\'erique.
		\emph{Comptes Rendus Acad. Sciences Paris}
		\textbf{222} (\textbf{1946}), 847--849.

		\bibitem{singh2007} Singh, G., M\'emoli, F., Carlsson, G. Topological
		Methods for the Analysis of High Dimensional Data Sets and 3D Object
		Recognition. In \emph{PBG@ Eurographics} \textbf{2007}, pp. 91--100.

		\bibitem{lum2013} Lum, P. Y., Singh, G., Lehman, A., Ishkanov, T.,
		Vejdemo-Johansson, M., Alagappan, M., Carlsson, J., Carlsson, G.
		Extracting insights from the shape of complex data using topology.
		\emph{Scientific Reports} \textbf{3} (\textbf{2013}), 1236.

		\bibitem{carriere2018} Carri\`ere, M., Oudot, S. Structure and
		Stability of the One-Dimensional Mapper. \emph{Foundations of
		Computational Mathematics} \textbf{18} (\textbf{2018}), 1333--1396.

	\end{thebibliography}

	\bigskip
	\begin{flushleft}
		(Guilherme Vituri)\,\, \textsc{[Simbolic Mind]}\\
		\emph{Email address:}  \texttt{viturivituri@gmail.com}
	\end{flushleft}

\end{document}
